%! Author = zhiweizhang
%! Date = 10.06.25

\chapter{Literature Review}\label{chapter:literature-review}

The classification of Nazi symbols in images is a highly specialized task situated at the intersection of computer
vision, hate speech detection, few-shot learning, and ethical \ac{AI} deployment. Owing to the diverse and sensitive
visual characteristics of such symbols, this problem introduces unique challenges, including the need for robust
visual representations, limited annotated data, and open-set recognition capabilities.

\section{Visual Symbol and Logo Recognition}

Foundational work in logo and symbol recognition has informed methods for visual hate symbol detection. Iandola et
al.~\parencite{iandola2015deeplogo} demonstrated that \ac{CNN}s are effective at recognizing logos in natural
scenes, which closely parallels the challenges of detecting hate symbols embedded in visually complex or cluttered
environments. Expanding on scalability and intra-class variability, Wang et al.~\parencite{li2020logo2k} proposed
Logo-2K+, a large-scale logo dataset comprising over 2,000 categories, highlighting the importance of generalization
across visually similar but distinct symbol classes.

Real-world applications of Nazi symbol recognition must also contend with stylistic variations, modifications, and
unseen symbol variants. Tüzkö et al.~\parencite{tuzko2020open} addressed these issues through open-set logo
detection techniques, enabling models to differentiate between known and unknown logos. This capability is essential
for the classification of hate symbols, which often appear in modified or contextually obscured forms.

\section{Hate Content Detection in Social Media and Memes}

Hate symbol recognition increasingly intersects with multimodal content analysis, especially in memes and online
extremist media. Belete and Alitasb~\parencite{belete2025200258} developed a deelp learning pipeline for detecting
offensive Amharic-language memes using \ac{BiLSTM} and dense networks. Although language-specific, their work
demonstrates the potential of combining textual and visual cues for hate content detection—a key consideration for
memes that integrate Nazi symbols with provocative slogans or captions.

Das Bhattacharjee et al.~\parencite{dasbhattacharjee2019365} adopted a context-aware active learning framework to
identify malicious content on social platforms. Their results underscore the importance of contextual
information—particularly relevant when hate symbols are embedded in complex visual narratives or are visually
subtle. In parallel, Scrivens et al.~\parencite{scrivens2021dangerous} examined how violent right-wing extremist
groups use imagery in online spaces, offering sociological insights that reinforce the need for comprehensive visual
detection approaches capable of capturing both overt and covert symbolism.

\section{Deep Learning Models for Visual Classification}

Modern image classification systems largely rely on deep neural networks. AlexNet~\parencite{krizhevsky2012imagenet}
marked a significant milestone in \ac{DL}, introducing a convolutional framework that enabled large-scale
visual classification. Subsequent advancements such as \ac{VGGNet}~\parencite{simonyan2014very} demonstrated that
deeper networks can capture finer-grained visual patterns—an important feature when symbols are small, degraded, or
partially obscured.

More recently, the \ac{ViT} architecture introduced by Dosovitskiy et al.~\parencite{dosovitskiy2020image} replaced
convolutions with self-attention mechanisms, enabling models to process
global image context. This is particularly beneficial for recognizing Nazi symbols that may appear in cluttered or
visually dense settings, such as posters, flags, or memes.

\section{Zero-Shot and Few-Shot Symbol Recognition}

A persistent challenge in Nazi symbol classification is the scarcity of labeled training data, especially for less
common or region-specific variations. In such cases, zero-shot and few-shot learning methods become essential. Xian
et al.~\parencite{xian2017zeroshot} systematically benchmarked \ac{ZSL} approaches,
underscoring the trade-offs between generalization and domain adaptation.

\ac{CLIP}, introduced by Radford et al.~\parencite{radford2021clip}, leverages contrastive learning to align image and
text representations, allowing for zero-shot classification based on descriptive prompts. This modality is
especially promising for detecting evolving or rarely seen Nazi symbols. DiffCLIP~\parencite{zhang2021language}
extends this work by enabling prompt-based fine-tuning, making it suitable for few-shot scenarios where only minimal
labeled data are available.

\ac{RN}~\parencite{sung2018learning} offer an alternative few-shot approach by classifying
unseen categories based on feature similarity. These methods are highly relevant in contexts where limited annotated
datasets of hate symbols must be leveraged for generalization to new symbol types or stylized variants.

\section{Dataset Transparency and Ethical Considerations}

The sensitive nature of Nazi symbol detection necessitates responsible and transparent \ac{AI} practices. Gebru et
al.~\parencite{gebru2021datasheets} introduced the concept of “datasheets for datasets,” advocating for
comprehensive documentation to promote dataset transparency, provenance tracking, and bias mitigation. Such
practices are vital when working with hate-related or politically sensitive imagery.

Additionally, Shahriari~\parencite{ieee2021ethics} emphasized ethically aligned design in \ac{AI} systems,
encouraging practitioners to prioritize human rights, safety, and fairness. These principles are particularly
salient for Nazi symbol classifiers, where the risk of misuse, false positives, or sociopolitical harm must be
carefully mitigated through responsible design and deployment.

